\documentclass[a4paper]{article}
% Español (México) con XeLaTeX: fontspec para fuentes Unicode, polyglossia
% para la división silábica y los nombres del documento en la variante mexicana.
\usepackage{fontspec}
\usepackage{polyglossia}
\setdefaultlanguage[variant=mexican]{spanish}
% Los nombres chinos van como «pinyin (original)»: xeCJK da a los caracteres
% Han su propia fuente; sin ella XeLaTeX los omite con un aviso.
% FandolSong no cubre algunos caracteres de nombres (U+73FA); WenQuanYi Zen Hei sí.
\usepackage[AutoFallBack=true]{xeCJK}
\setCJKmainfont{FandolSong-Regular.otf}
\setCJKfallbackfamilyfont{\CJKrmdefault}{WenQuanYi Zen Hei}
% polyglossia no traduce los nombres de listings.
\AtBeginDocument{\providecommand{\lstlistingname}{}\renewcommand{\lstlistingname}{Listado}%
  \providecommand{\lstlistlistingname}{}\renewcommand{\lstlistlistingname}{Índice de listados}}
\usepackage{amsmath, amssymb}
\usepackage{graphicx}
\usepackage[margin=2.5cm]{geometry}
\usepackage[most]{tcolorbox}
\usepackage{etoolbox}
\usepackage{listings}
\usepackage{hyperref}
\usepackage{booktabs}
\usepackage{subcaption}
\usepackage{float}
\usepackage{array}
\newtcolorbox{knowledgebox}[1]{enhanced, colback=blue!5!white, colframe=blue!75!black, colbacktitle=blue!75!black, coltitle=white, fonttitle=\bfseries, title={#1}, attach boxed title to top left={yshift=-2mm, xshift=2mm}, boxrule=1pt, sharp corners}
\newtcolorbox{importantbox}[1]{enhanced, colback=yellow!10!white, colframe=yellow!80!black, colbacktitle=yellow!80!black, coltitle=black, fonttitle=\bfseries, title={#1}, sharp corners}
\newtcolorbox{warningbox}[1]{enhanced, colback=red!5!white, colframe=red!75!black, colbacktitle=red!75!black, coltitle=white, fonttitle=\bfseries, title={#1}, sharp corners}
\lstset{language=Python, basicstyle=\ttfamily\small, keywordstyle=\color{blue}, stringstyle=\color{red!60!black}, commentstyle=\color{green!60!black}, breaklines=true, frame=single, numbers=left, numberstyle=\tiny\color{gray}, captionpos=b, extendedchars=true}
\newcommand{\notetitle}{Carnaval de IA de Qingke (青稞): discusión temática sobre Infra}
\newcommand{\noteauthors}{Elaborado a partir de la discusión pública del Carnaval de IA de Qingke}
\newcommand{\notedate}{2026-04-02}
\newcommand{\videochannel}{Comunidad Qingke}
\newcommand{\videopublishdate}{2025}
\newcommand{\videoduration}{aproximadamente 112 minutos}
\newcommand{\videourl}{}
\newcommand{\videocoverpath}{cover.jpg}
\newcommand{\repourl}{https://github.com/hqhq1025/ai-course-notes}

\usepackage{fancyhdr}
\pagestyle{fancy}
\fancyhf{}
\fancyfoot[C]{\thepage}
\fancyfoot[R]{\footnotesize\color{gray}\href{\repourl}{AI Course Notes}}
\renewcommand{\headrulewidth}{0pt}
\renewcommand{\footrulewidth}{0pt}

\begin{document}
\begin{titlepage}
\centering
{\Large Notas de entrevista\par}
\vspace{1.2cm}
{\huge\bfseries \notetitle\par}
\vspace{0.8cm}
{\large \noteauthors\par}
\vspace{0.3cm}
{\large \notedate\par}
\vspace{1.2cm}
\ifdefempty{\videocoverpath}{}{\includegraphics[width=0.82\textwidth,height=0.45\textheight,keepaspectratio]{\videocoverpath}\par}
\vfill
\begin{tcolorbox}[width=0.9\textwidth, colback=black!2!white, colframe=black!60, sharp corners]
\textbf{Autor o canal del video}: \videochannel\par
\textbf{Fecha de publicación}: \videopublishdate\par
\textbf{Duración del video}: \videoduration\par
\ifdefempty{\videourl}{}{\textbf{Enlace del video}: \href{\videourl}{\nolinkurl{\videourl}}\par}
\textbf{Página del proyecto}: \href{\repourl}{\nolinkurl{\repourl}}\par
\end{tcolorbox}
\vspace{0.3cm}
{\small Esta charla se concentra en dos desplazamientos de enfoque en la infraestructura de los modelos grandes durante 2025: en la primera mitad del año, en torno a la clusterización de la inferencia y la separación PD; en la segunda mitad, en torno al scaling de la infraestructura de RL / Agent. Los invitados provienen tanto del ámbito académico como de la industria, por lo que la charla hace énfasis particular en «cómo el sistema absorbe nuevas workloads».\par}
\end{titlepage}
\tableofcontents
\newpage

\section{El tema central de la infraestructura en 2025: de la clusterización de la inferencia a la infraestructura de RL / Agent}

\subsection{Por qué la infraestructura ocupó por primera vez el centro del escenario}
La apertura del panel planteó el problema con mucha claridad: en 2025 la infraestructura ya no es solo un soporte de fondo, sino la condición previa para que la capacidad de un modelo pueda materializarse. Varios de los invitados mencionaron una sensación compartida: en el pasado, la discusión de la industria se concentraba más en el modelo mismo, la innovación algorítmica y la escala de los datos, pero hacia 2025 empezó a repetirse la conciencia de que ``el mismo modelo, si la infraestructura no funciona, no puede liberar su capacidad bajo una carga real''.

\begin{importantbox}{Por qué la importancia de la infraestructura aumentó de manera notoria en 2025}
Hay al menos tres razones:
\begin{itemize}
  \item la inferencia pasó de experimentos en una sola máquina a producción a gran escala, por lo que el rendimiento, la latencia y la estratificación de recursos deben diseñarse de forma sistemática;
  \item RL, Agent y multi-modality trajeron consigo rollouts más largos, workloads más heterogéneos y cadenas de retroalimentación más complejas;
  \item la industria ya no solo entrena modelos, sino que empieza a entrenar todo un sistema de servicio capaz de interactuar con sistemas externos.
\end{itemize}
\end{importantbox}

Este cambio significa que la infraestructura ya no consiste solo en ``hacer que los operadores corran más rápido'', sino que empieza a determinar el paradigma mismo de la investigación. Por ejemplo, cuando se discute si Agent RL es viable, en el fondo se discute si el motor de rollout, la construcción del entorno, la combinación de entrenamiento e inferencia y la recuperación ante fallas son suficientes para sostener el ciclo cerrado del experimento.

\subsection{La palabra clave del primer semestre: la clusterización de la inferencia}
Zhang Mingxing (张明星) ofreció una división muy nítida: el tema del primer semestre de 2025 fue la clusterización de la inferencia, mientras que el del segundo semestre se inclinó más hacia el scaling de la infraestructura de RL. Esta división resulta muy reveladora, porque muestra que la industria no resuelve todos los problemas al mismo tiempo, sino que primero convierte la inferencia en línea en una base industrializada, y solo después conecta a ella workloads más complejos como RL y Agent.

\textit{``En el primer semestre, en realidad todo se trataba de la clusterización de la inferencia; el tema del segundo semestre fue, básicamente, cómo hacer scaling up / scaling out de la infraestructura de RL.''}

\begin{knowledgebox}{Por qué la clusterización de la inferencia es condición previa de la infraestructura de RL / Agent}
Porque el cuello de botella de muchos workloads nuevos no está en el backprop, sino en el rollout, el sampling y la interacción con el entorno. Dicho de otro modo, si el motor de inferencia no puede funcionar de manera estable, con alto rendimiento y baja fragmentación, entonces el entrenamiento posterior de RL / Agent simplemente no cuenta con un ciclo de retroalimentación suficientemente barato.
\end{knowledgebox}

\subsection{Resumen de la sección}
La infraestructura fue importante en 2025 no porque la ingeniería de sistemas se haya vuelto de repente un tema de moda, sino porque tanto el entrenamiento de modelos como el despliegue de productos se vieron forzados a una etapa en la que ``deben sistematizarse''. La clusterización de la inferencia se convirtió primero en el eje del primer semestre, y después la infraestructura de RL / Agent tomó el relevo, conformando el eje central de la discusión de todo el año.
\section{La madurez de la infraestructura de inferencia: separación PD, EP grande y colaboración nube-borde}

\subsection{Por qué la separación PD se popularizó tan rápido}
La primera palabra clave que el invitado mencionó repetidamente fue la separación PD, es decir, Prefill-Decode Disaggregation. La razón por la que en 2025 pasó rápidamente de ser un concepto de investigación a una arquitectura de producción es que las dos etapas del workload de inferencia tienen necesidades de hardware completamente distintas: Prefill consume más cómputo, mientras que Decode consume más ancho de banda y gestión del KV cache. Atar ambas etapas a un mismo conjunto de máquinas equivale, en esencia, a desperdiciar recursos.

\begin{importantbox}{El valor de ingeniería de la separación PD}
La separación PD permitió, por primera vez, que los clústeres de inferencia contaran con una forma de orquestación de recursos con mayor granularidad:
\begin{itemize}
  \item los nodos de Prefill pueden optimizarse para un alto poder de cómputo;
  \item los nodos de Decode pueden optimizarse para un ancho de banda alto y una gran capacidad de caché;
  \item el sistema en su conjunto puede escalar de forma independiente con mayor facilidad según los cambios de tráfico;
  \item la velocidad de transferencia de tecnología entre la comunidad de código abierto y los despliegues industriales se aceleró notablemente.
\end{itemize}
\end{importantbox}

Zhang Mingxing (张明星) mencionó que, cuando apenas comenzó el proyecto relacionado el año pasado, todavía había pocos usuarios; para 2025, impulsados por la colaboración de la comunidad, casi todos los proveedores principales habían adoptado arquitecturas similares. Esto ilustra bien una tendencia: cuando el workload se vuelve suficientemente grande, la optimización a nivel de arquitectura llega a un consenso en la industria más rápido que la optimización local de kernels.

\subsection{EP grande e inferencia con MoE: otra línea de expansión de capacidad}
Otra línea principal de la infraestructura de inferencia es el Expert Parallelism a gran escala. A medida que los modelos MoE siguen expandiéndose, el costo del enrutamiento de expertos y de la comunicación entre tarjetas es cada vez mayor; la Infra ya no solo tiene que responder «si el modelo cabe o no», sino también «cómo dividir a los expertos», «cómo evitar que la comunicación entre nodos arrastre la latencia» y «cómo equilibrar el servicio entre el rendimiento y la latencia de cola».

\begin{warningbox}{La dificultad de la Infra de MoE no es simplemente añadir más tarjetas}
Cuantos más Experts hay, más complejos se vuelven el enrutamiento, el balanceo de carga y la transferencia de activaciones entre dispositivos. Si la estrategia de programación y la topología de despliegue no están bien diseñadas, un EP grande puede, por el contrario, hacer que el sistema salga perdiendo en comunicación y fragmentación.
\end{warningbox}

\subsection{El lado de la nube y el lado del borde no son el mismo tipo de problema de optimización}
Otro punto práctico del panel es distinguir con claridad entre la inferencia en la nube y la inferencia en el borde. Los dos proyectos de código abierto en los que trabaja el invitado se orientan, uno hacia la optimización de la inferencia en clústeres distribuidos del lado de la nube, y el otro hacia la optimización de la inferencia heterogénea CPU-GPU del lado del borde. Esto corresponde a dos objetivos de sistema completamente distintos: el lado de la nube busca el rendimiento total, la agrupación de recursos y la eficiencia multiinquilino; el lado del borde se preocupa más por la latencia de una sola máquina, el consumo de energía y la privacidad local.

\begin{table}[H]
\centering
\begin{tabular}{p{2.8cm}p{4.0cm}p{4.5cm}}
\toprule
\textbf{Forma de despliegue} & \textbf{Objetivo principal} & \textbf{Puntos típicos de atención de la Infra} \\
\midrule
Clúster en la nube & Rendimiento, escalado, utilización multiinquilino & Separación PD, EP grande, comunicación entre nodos, orquestación de servicios \\
Borde & Baja latencia, bajo consumo de energía, privacidad, disponibilidad local & Heterogeneidad CPU-GPU, poda de operadores, uso de memoria, adaptación de dispositivos \\
Modo híbrido & Equilibrio entre costo y experiencia & Localización de la ruta caliente, migración de la ruta fría a la nube, sincronización de estado \\
\bottomrule
\end{tabular}
\caption{Aunque ambos se llaman «optimización de inferencia», los objetivos de optimización del lado de la nube y del lado del borde no son los mismos}
\end{table}

\subsection{Resumen de la sección}
El avance más importante de la Infra de inferencia en 2025 no es el nombre de un proyecto en particular, sino que toda la industria llegó a un consenso más claro sobre la optimización a nivel de arquitectura: la separación PD, el EP grande y la colaboración nube-borde nos indican que la inferencia ya evolucionó de «hacer que funcione en una sola tarjeta» a «reconstruir el sistema según las características del workload».
\section{RL / Agent Infra: cuando el Actor crece, el Rollout se alarga y el entorno se vuelve más complejo}

\subsection{El cuello de botella de RL cada vez se parece más a un problema de inferencia}
Varios panelistas señalaron explícitamente que hoy el principal punto de bloqueo de la infraestructura de RL es el rollout, no el entrenamiento en sí. Esto es clave, porque explica por qué el motor de inferencia, la gestión de caché, las interfaces de entorno y el flujo de retorno de datos vuelven a ocupar el centro del diseño de sistemas en la era de RL. Sobre todo cuando el modelo crece y el rollout se alarga, los marcos de RL tradicionales se enfrentan a un Actor completamente distinto del de antes.

\begin{importantbox}{La diferencia esencial entre el RL de los grandes modelos actuales y el RL tradicional}
Los panelistas compararon esto con los sistemas de MuJoCo / Atari de años anteriores y señalaron que la diferencia se concentra en tres puntos:
\begin{itemize}
  \item el Actor pasó de una red de unos pocos millones de parámetros a un enorme LM / VLM;
  \item el entorno pasó de un simulator pequeño y relativamente fijo a herramientas y servicios reales, más dinámicos y heterogéneos;
  \item el rollout ya no es un forward barato, sino un proceso completo de inferencia e interacción.
\end{itemize}
\end{importantbox}

\subsection{El rollout largo y el modelo grande multiplican la presión sobre el sistema}
Los panelistas mencionaron varios retos muy concretos: cómo hacer RL cuando el modelo es más grande, cómo hacer RL cuando el rollout es especialmente largo, cómo hacer RL con multi-modality y con Agent workload. Estos problemas parecen algorítmicos, pero en realidad se traducen directamente en presión sobre el sistema. Por ejemplo, que el rollout se alargue implica mantener la caché durante más tiempo, una planificación más compleja y un mayor costo de recuperación ante fallos; que el entorno sea más complejo implica que el marco de entrenamiento debe soportar flujos de datos y de estado más heterogéneos.

\begin{knowledgebox}{Por qué la infraestructura de RL necesita integrar entrenamiento e inferencia}
Si el marco de entrenamiento y el marco de inferencia están completamente separados, surgen varios problemas:
\begin{itemize}
  \item la cadena de retorno de los resultados del rollout se vuelve demasiado larga;
  \item la gestión del estado queda dispersa en múltiples sistemas;
  \item al depurar una falla, resulta difícil determinar si está en el entrenamiento, en la inferencia o en el entorno;
  \item el lado algorítmico y el lado de sistemas suelen ``mirarse el uno al otro sin entenderse'', cada uno cree que el cuello de botella está del lado contrario.
\end{itemize}
\end{knowledgebox}

\subsection{El RL de Agent pone en primer plano el problema de construir el entorno}
El RL tradicional suele suponer que el entorno ya está dado, pero en las tareas de Agent, el entorno mismo forma parte del diseño del sistema. En el panel se mencionó varias veces environment construction, workflow construction, real world task, failure handling, entre otros temas, lo cual muestra que el RL de Agent no consiste solo en ``aprender una política dentro de un entorno ya existente'', sino que con frecuencia hay que construir primero el entorno antes de que el modelo pueda aprender en él.

\begin{warningbox}{Si el entorno no está bien construido, el RL aprende de manera sesgada}
Los panelistas mencionaron un punto de dolor central: muchos sistemas se entrenan muy bien dentro de un solo entorno, pero en cuanto se cambia de entorno dejan de generalizar. Es decir, las diferencias en la construcción del entorno, en la estructura de la tarea y en la forma de la recompensa se convierten directamente en un problema de generalización. Este no es un error que se pueda eliminar simplemente aumentando la cantidad de datos.
\end{warningbox}

\subsection{Resumen de la sección}
La complejidad de la infraestructura de RL / Agent proviene de que tres variables cambian al mismo tiempo: el Actor crece, el rollout se alarga y el entorno se vuelve más complejo. Por eso el centro del diseño de sistemas también se desplaza de ``lograr que el entrenamiento funcione'' a ``lograr que el entrenamiento, la inferencia, el entorno y el flujo de retorno de datos formen un ciclo reproducible y escalable''.
\section{El sistema de entrenamiento está pasando de un marco de línea principal única a un sistema de programación heterogénea}

\subsection{Por qué los marcos de entrenamiento tradicionales ya no bastan}
En el panel hay una observación muy precisa: los marcos optimizados en los últimos dos años en torno al pre-training son, en esencia, un caso especial ``demasiado apegado a una línea principal''. La longitud de la entrada, la estructura de la pasada hacia adelante y el flujo de datos eran relativamente uniformes, de modo que el sistema podía optimizarse con facilidad en torno a una sola línea principal. Pero al llegar a RL, Agent, multimodal y workload impulsado por workflow, la estructura de entradas y salidas se vuelve mucho más heterogénea, y muchos supuestos del sistema comienzan a fallar.

\textit{``Este año, al diseñar el sistema de entrenamiento, el énfasis está en cómo modularizar y separar los problemas de infra, y al mismo tiempo unir el marco de inferencia con el marco de entrenamiento.''}

\begin{importantbox}{El nuevo objetivo del sistema de entrenamiento no es ser más único, sino más modular}
La nueva dirección de diseño se parece más a esto:
\begin{itemize}
  \item que el marco de entrenamiento y el marco de inferencia compartan interfaces más estables;
  \item modularizar environment, reward, workflow y failure handling;
  \item permitir que distintos workload corran en distinto hardware y distintas configuraciones de precisión;
  \item que el sistema primero absorba la heterogeneidad, y solo después se hable de exprimir el rendimiento.
\end{itemize}
\end{importantbox}

\subsection{Las fallas, el sistema de archivos y el hardware heterogéneo se vuelven la norma}
Muchos de los detalles de los que hablaron los invitados tienen un sabor muy ``de Infra'': qué hacer cuando falla el montaje del sistema de archivos, cómo recuperarse cuando el entrenamiento tiene una falla, cómo manejar las diferencias de precisión y estabilidad entre distintas tarjetas, cómo programar el entrenamiento entre clusters. La señal que se desprende de esto es clara: el sistema de entrenamiento de la era de Agent / RL ya no opera en condiciones ideales de laboratorio, sino que corre por defecto en un mundo que falla, que fluctúa y que es heterogéneo.

\begin{knowledgebox}{Por qué se vuelven importantes los sistemas event-driven y dinámicos}
El workflow de un Agent ya no es un flujo de tensores de trayectoria única, sino algo que percibe el entorno, invoca herramientas, espera retroalimentación y luego decide el siguiente paso. Esto se aproxima de forma natural a un event-driven system, no a un batch training estático. Lo que el sistema necesita procesar ya no son solo tensores, sino eventos asíncronos, actualizaciones de estado y workflow heterogéneo.
\end{knowledgebox}

\subsection{El sistema y el algoritmo deben evolucionar juntos}
Varios invitados mencionaron un escenario frecuente: el lado del algoritmo dice que el sistema es demasiado lento, y el lado del sistema dice que el workload es demasiado irregular. Esta tensión no es algo negativo; más bien muestra que la era de Agent ya puso la ``innovación algorítmica'' y el ``soporte de sistema'' en la misma mesa. Antes se podía primero construir el algoritmo y después dejar que el sistema se adaptara; ahora, muchas capacidades, sin la Infra correspondiente, ni siquiera pueden validarse de forma confiable.

\begin{warningbox}{Perseguir solo el rendimiento del sistema, o solo la novedad del algoritmo, distorsiona el resultado}
Si el rendimiento del sistema es muy alto pero solo se adapta a un único workload, entonces deja de funcionar en cuanto cambia el entorno; si la idea algorítmica es muy novedosa pero el sistema no puede reproducirla de forma estable, tampoco puede convertirse en un consenso de la comunidad. Lo que en verdad enfatiza el tema de Infra es que ambos deben alcanzar la misma velocidad de iteración.
\end{warningbox}

\subsection{Resumen de la sección}
El sistema de entrenamiento está pasando de un ``marco de línea principal única'' a un ``sistema de programación heterogénea''. Detrás de este cambio está la transformación del workload: de un pre-training uniforme hacia un workload de RL / Agent que es event-driven, intensivo en environment y propenso a fallas. El centro del diseño de sistemas ha pasado del rendimiento estático a la capacidad de sostenerse de forma dinámica.
\section{Panorama 2026: no se entrena solo un modelo, sino todo un sistema}

\subsection{Del modelo como centro al sistema como centro}
Hacia el final del panel hay una frase que vale la pena recordar: el futuro no es solo entrenar modelos, sino entrenar un sistema. Aquí el sistema incluye el renderizado, la inferencia, los flujos de trabajo, la construcción del entorno, el retorno de datos y la interfaz de despliegue que rodean al modelo. Esto significa que el «modelo» ya no es el único objeto de optimización, sino apenas el eslabón más central, pero no el único, de toda la cadena.

\begin{importantbox}{Direcciones con alta probabilidad de seguir expandiéndose en 2026}
\begin{itemize}
  \item La combinación de inferencia y entrenamiento seguirá profundizándose; la separación entre entrenamiento e inferencia dejará de ser un problema de arquitectura para convertirse en un problema de workload;
  \item La construcción del Agent environment se convertirá en una capacidad de ingeniería importante;
  \item La multimodalidad, el RL y los workloads de Agent seguirán impulsando la heterogeneidad del sistema;
  \item La experiencia histórica de los sistemas distribuidos, las bases de datos y la tolerancia a fallos se reintroducirá en la Infra de los modelos grandes.
\end{itemize}
\end{importantbox}

\subsection{El destino final de la Infra no es esconderse, sino convertirse en capacidad por defecto}
Por el tono del panel se percibe que los invitados no buscan que el usuario perciba directamente la Infra, sino que esta termine retrocediendo hacia el fondo y se convierta en una capacidad de soporte que el sistema posee por defecto. Es decir, cuando discutimos si «un producto de Agent es fácil de usar», muchos de los factores que realmente lo deciden están escondidos en el motor de rollout, la estrategia de caché, el sistema de programación, la interfaz del entorno y la recuperación ante fallos.

\begin{table}[H]
\centering
\begin{tabular}{p{2.8cm}p{4.2cm}p{4.4cm}}
\toprule
\textbf{Dirección} & \textbf{Objetivo del sistema} & \textbf{Pregunta clave del lado de la Infra} \\
\midrule
La inferencia se sigue agrupando en clústeres & Mayor rendimiento, menor latencia de cola & Separación PD, paralelismo de expertos, programación entre nodos \\
Escalamiento de RL / Agent & Rollouts más largos, entornos más grandes, retroalimentación más compleja & Combinación de entrenamiento e inferencia, interfaz de environment, retorno de reward \\
Sistemas de entrenamiento heterogéneos & Soportar más combinaciones de workload y hardware & Abstracción modular, recuperación ante fallos, gestión de precisión \\
Adopción industrial & Que el modelo se conecte de verdad a la cadena de la aplicación & Sistema de archivos, permisos, workflow, estabilidad \\
\bottomrule
\end{tabular}
\caption{Toda la discusión resumida en «qué es lo siguiente que la Infra debe absorber»}
\end{table}

\subsection{Resumen de la sección}
Este panel dedicado a la Infra termina en un juicio muy claro: la siguiente etapa de los sistemas de modelos grandes no consiste simplemente en apilar más tarjetas, sino en hacer que la inferencia, el entrenamiento, el entorno y el flujo de trabajo formen un sistema integral verdaderamente escalable. Esta es también la razón de fondo por la que la Infra vuelve a convertirse en el eje central en la era de los Agent.
\section{Síntesis y ampliación}

\subsection{Panorama de las ideas centrales}
\begin{table}[H]
\centering
\begin{tabular}{p{2.8cm}p{4.0cm}p{5.0cm}}
\toprule
\textbf{Tema} & \textbf{Conclusión del panel} & \textbf{Problema de sistemas correspondiente} \\
\midrule
Infraestructura de inferencia & La separación PD y el EP grande ya son el estándar en producción & ¿Cómo reorganizar los recursos según la carga de trabajo? \\
Infraestructura de RL & El rollout es el cuello de botella central; combinar entrenamiento e inferencia es cada vez más importante & ¿Cómo hacer que el RL de modelos grandes sea escalable? \\
Infraestructura de agentes & El entorno, el workflow y los flujos de eventos hacen el sistema más dinámico & ¿Cómo dar cabida a herramientas reales y fallos reales? \\
Sistema de entrenamiento & Se pasa de un framework de una sola línea principal a un sistema de programación heterogéneo & ¿Cómo separar en módulos los problemas de infraestructura? \\
Dirección futura & No se trata solo de entrenar el modelo, sino de entrenar el sistema & ¿Cómo hacer que el sistema y el algoritmo iteren juntos? \\
\bottomrule
\end{tabular}
\caption{Las cinco líneas más importantes que recordar del panel de infraestructura}
\end{table}

\subsection{Tres conclusiones orientadas a la práctica}
\begin{enumerate}
  \item Si la carga de trabajo ya se volvió Agent / RL / multimodal, no se debe seguir diseñando el sistema con la mirada de un preentrenamiento puro.
  \item Lo que cada vez más limita el techo de las capacidades de los modelos grandes es el rollout, el entorno, la programación y la recuperación de fallos, no la velocidad de un solo operador.
  \item La señal de que una infraestructura es madura no es que el usuario perciba su complejidad, sino que se mantenga estable, reproducible y escalable bajo cargas de trabajo complejas.
\end{enumerate}

\begin{importantbox}{La verdadera enseñanza del panel de infraestructura}
Nos recuerda que, en la era de los modelos grandes, el ``sistema'' no es un accesorio, sino parte de la capacidad. Muchos problemas que parecen cuellos de botella algorítmicos solo pueden resolverse de verdad cuando el diseño del sistema, la construcción del entorno y las interfaces de despliegue avanzan al mismo ritmo.
\end{importantbox}

\subsection{Lecturas adicionales}
\begin{itemize}
  \item Material sobre sistemas de inferencia como MoonCake, KTransformers, vLLM y SGLang
  \item Proyectos de código abierto y artículos relacionados con la infraestructura de RL y los rollout engines
  \item Trabajos sobre entornos de agentes, modelos de mundo y construcción de benchmarks
  \item Experiencia de migración de entrenamiento distribuido, recuperación ante fallos, sistemas de archivos y bases de datos hacia la infraestructura de IA
  \item Prácticas de ingeniería sobre heterogeneous workflow y sistemas de IA basados en eventos
\end{itemize}

\end{document}
